\documentclass[11pt,a4paper]{article}
\usepackage[T1]{fontenc}
\usepackage{isabelle,isabellesym}
\usepackage{pdfsetup}
\hypersetup{hidelinks}

\urlstyle{rm}
\isabellestyle{it}

\begin{document}

\title{Set-Coded Computation in Isabelle/ZF}
\author{Tang Ziyi}
\maketitle

\begin{abstract}
We formalise deterministic binary Turing machines as sets in Isabelle/ZF.
Finite instruction tables determine execution on a tape represented by two
finite lists. We prove that equality of observed tape contents is preserved
by execution, construct natural-number codes for machines and configurations,
and prove a numerical evaluator correct for any finite number of steps.
We certify its numerical operations and finite-step evaluation as
primitive recursive.
A concrete finite rejection transformation yields the undecidability of
self-input halting in this machine model. We also prove that an explicitly
supplied blank buffer protects saved tape data through every prefix of a
bounded execution. Numeric evaluation is a mathematical function on codes;
the entry does not construct a universal Turing machine.
\end{abstract}

\tableofcontents

\section{Scope and representation}
The machine instruction format follows the five-action operational design of
the Isabelle/HOL AFP entry \emph{Universal Turing Machine}~\cite{afp-utm}.
Here machines, configurations, and finite runs are sets in Isabelle/ZF.
An instruction separately writes a symbol, moves the head, or leaves the tape
unchanged, then selects the next control state. The final state is absorbing;
an absent instruction enters it without changing the tape.

The two finite tape lists represent a blank-ended tape. Literal list equality
distinguishes representations that have the same observed contents, such as a
list ending in one blank and the same list without that blank. The entry
therefore defines equality by the symbols observed at every position relative
to the head and proves that execution respects it. This is also necessary for
the numerical output interpretation in the separate research development.

The machine numbering and numerical evaluator are proved correct as set
functions. Primitive-recursive certificates for operations on codes and
bounded evaluation use the Isabelle/ZF theory ZF-Induct.Primrec. They do not
construct a Turing machine implementing the evaluator. The undecidability
proof uses a different branch of the development. It supplies the finite
rejection transformation required by the diagonal argument, rather than
assuming that transformation exists. The workspace theorem is a third
branch. It protects data appended outside an explicitly bounded region
without predicting the duration of a program. Neither the evaluator nor
the workspace theorem is used by the diagonal proof.

Inputs are unary lists of ones. A halted machine accepts or rejects according
to the symbol under its head. This output convention belongs to the model
used here and is not inherited from the Isabelle/HOL entry cited above.

The formalisation is motivated in part by the classical question of machine
computability and its limits~\cite{turing1937}. The results below concern the
machine model defined here. No equivalence theorem with another machine
formalisation, general compiler, or universal-machine theorem is asserted.
Set-coded finite runs can serve as parameters in later set-theoretic model
developments. No internal halting formula or satisfaction theorem is proved
in this entry.

\input{session}

\section*{Development provenance}
Tang Ziyi formulated the research question and wrote the initial Isabelle
formalisation. Subsequent definitions, theorem statements, proofs, and
repository engineering include substantial assistance from OpenAI Codex.

\begin{thebibliography}{9}

\bibitem{afp-utm}
Jian Xu, Xingyuan Zhang, Christian Urban, Sebastiaan J. C. Joosten, and
Franz Regensburger.
\emph{Universal Turing Machine}.
Archive of Formal Proofs, 2019, revised 2022.
\url{https://isa-afp.org/entries/Universal_Turing_Machine.html}.

\bibitem{turing1937}
Alan M. Turing.
On computable numbers, with an application to the Entscheidungsproblem.
\emph{Proceedings of the London Mathematical Society}, series 2, volume 42,
pages 230--265, 1937.
\url{https://doi.org/10.1112/plms/s2-42.1.230}.

\end{thebibliography}

\end{document}
